\chapter{What a chess engine is actually trying to do}
\label{ch:problem}

\begin{objectives}
\begin{itemize}
\item State the engine problem as exactly two questions, and show that every engine ever
  built is a different answer to those two questions.
\item Compute, from your own 9{,}030-game corpus, why brute force is not merely slow but
  physically impossible.
\item Run minimax and alpha--beta by hand on a small tree, and count the nodes each visits.
\item Write down the classical evaluation function $f(s)=\sum_i w_i \phi_i(s)$ and name
  precisely the two things it cannot do.
\item Explain why Leela replaces \emph{both} halves of the classical answer, not just one.
\end{itemize}
\end{objectives}

\section{Two questions, and only two}

Put a position in front of any chess-playing machine ever built --- Turing's paper machine
of 1948, Deep Blue, Stockfish, Leela --- and it must answer two questions.

\begin{keyidea}
\textbf{Question 1 (evaluation).} \emph{How good is this position?} Reduce a board to a
number, or to something a number can be extracted from.

\medskip
\textbf{Question 2 (search).} \emph{Where should I spend my thinking?} There are more
continuations than can ever be examined. Which ones get examined?
\end{keyidea}

That is the whole subject. Everything in the rest of this book --- transformers, attention
heads, Monte-Carlo trees, linear probes --- is machinery bolted onto one of those two
questions. It is worth fixing the shape of the problem in your head now, because the
distinctive thing about Leela is easy to state once you have it:

\begin{center}
\begin{tabular}{@{}lll@{}}
\toprule
 & \textbf{Answer to Q1 (evaluate)} & \textbf{Answer to Q2 (where to look)} \\
\midrule
Deep Blue, Stockfish (classical) & hand-written formula & alpha--beta, hand-written ordering \\
Stockfish (modern, NNUE)         & small neural network & alpha--beta, hand-written ordering \\
\textbf{Leela / AlphaZero}       & \textbf{neural network} & \textbf{neural network} + Monte-Carlo tree \\
\bottomrule
\end{tabular}
\end{center}

Leela's real novelty is the right-hand column. Plenty of engines learned their evaluation.
Leela is the family that also \emph{learned where to look} --- and that single fact is what
makes it interpretable in a way Stockfish is not, and therefore what makes your whole
project possible. When you ask Leela "what were you considering?", it has an answer, because
considering is something it does with a learned opinion rather than with a fixed procedure.

\section{Why you cannot simply look at everything}

The naive answer to Question~2 is: look at everything. Let us kill it with real numbers ---
yours.

\subsection{The branching factor, measured on your own games}

The \textbf{branching factor} $b$ is how many legal moves there are, on average, in a
position. Textbooks say "about 30" and rarely say where the number came from. Here it is,
measured on the corpus in \texttt{games\_of\_derdiedasdie/}:

\begin{realdata}[title={Real data: your corpus (\texttt{tools/corpus\_stats.py})}]
\begin{center}
\begin{tabular}{@{}lr@{}}
\toprule
Games & 9{,}030 \\
Positions examined & 549{,}816 \\
\textbf{Mean legal moves per position} & \textbf{31.35} \\
Median legal moves & 33 \\
Most legal moves seen in any position & 111 \\
Fewest (forced move) & 1 \\
Mean game length & 60.9 plies \\
\midrule
Mean legal moves, opening (moves 1--12) & 32.84 \\
Mean legal moves, middlegame & 32.48 \\
Mean legal moves, endgame ($\le 12$ pieces) & 15.81 \\
\bottomrule
\end{tabular}
\end{center}
\end{realdata}

So $b \approx 31$ in your chess. Note the shape of that table: the branching factor is
essentially flat at $\approx 32$ through opening and middlegame and then halves in the
endgame. That is one reason engines feel "deeper" in endgames --- the tree is genuinely
narrower there, so the same node budget buys more plies.

\begin{notationbox}
A \textbf{ply} is one move by one side. "White plays $e4$" is one ply. A \emph{move} in
chess notation ($1.e4\ e5$) is two plies. Engine depth is always counted in plies. In this
book, $d$ is depth in plies.
\end{notationbox}

\subsection{The count}

If every position has $b$ replies, then looking $d$ plies ahead exhaustively means visiting
about
\begin{equation}
  \label{eq:treesize}
  T(d) \;=\; b^{d}
\end{equation}
leaf positions. Put $b = 31$ in and turn the crank:

\begin{center}
\begin{tabular}{@{}rlr@{}}
\toprule
Depth $d$ (plies) & In chess terms & Positions $31^d$ \\
\midrule
2  & one move each        & $961$ \\
4  & two moves each       & $\approx 9.2\times10^{5}$ \\
6  & three moves each     & $\approx 8.9\times10^{8}$ \\
8  & four moves each      & $\approx 8.5\times10^{11}$ \\
10 & five moves each      & $\approx 8.2\times10^{14}$ \\
20 & ten moves each       & $\approx 6.7\times10^{29}$ \\
40 & a short game         & $\approx 4.5\times10^{59}$ \\
\bottomrule
\end{tabular}
\end{center}

\begin{worked}{how long would depth 20 take?}
Suppose you evaluate a billion ($10^9$) positions per second --- generous; a strong
classical engine on good hardware manages a few tens of millions.
\[
  \frac{6.7\times 10^{29}\ \text{positions}}{10^{9}\ \text{positions/second}}
  \;=\; 6.7\times 10^{20}\ \text{seconds}.
\]
A year is about $3.15\times10^{7}$ seconds, so this is
\[
  \frac{6.7\times10^{20}}{3.15\times10^{7}} \;\approx\; 2.1\times10^{13}\ \text{years},
\]
roughly \emph{fifteen hundred times the age of the universe} --- to see ten moves ahead
once. And your games last thirty moves, not ten.
\end{worked}

Shannon made this argument in 1950 and put the number of plausible chess games at about
$10^{120}$, now called the Shannon number. The exact exponent does not matter. What matters
is the shape of Equation~\eqref{eq:treesize}: the cost is \emph{exponential} in depth, so
hardware improvements buy almost nothing. A machine a thousand times faster buys you
\[
  \log_{31}(1000) \;\approx\; 2.0 \quad\text{extra plies.}
\]
One extra move each. That is what a factor of a thousand is worth. Every serious idea in
engine design is therefore an idea about \emph{not looking at things}.

\begin{keyidea}
Exponential growth means the interesting question is never "how fast can I evaluate?" but
"\textbf{how small can I make the effective branching factor?}" Write $b_{\text{eff}}$ for
the average number of moves an engine actually explores per node. Classical alpha--beta
with good move ordering gets $b_{\text{eff}}$ down to roughly $\sqrt{b} \approx 6$. Leela's
answer --- ask a network which moves are worth the visit --- is a different and, in the
positions this project cares about, more selective attack on the same quantity.
\end{keyidea}

\section{Minimax: the correct answer, if you could afford it}

Before improving on brute force, be precise about what brute force computes.

Chess is a two-player game with strictly opposed interests. If I score a position from
White's point of view, then White wants the score high and Black wants it low. Suppose we
can evaluate leaf positions somehow (that is Question~1; assume it for now). Then the value
of an internal node is defined recursively:

\begin{equation}
  \label{eq:minimax}
  V(s) \;=\;
  \begin{cases}
    f(s), & \text{if $s$ is a leaf (we stop here)},\\[4pt]
    \max_{a} V(s\!\cdot\!a), & \text{if White is to move at $s$},\\[4pt]
    \min_{a} V(s\!\cdot\!a), & \text{if Black is to move at $s$},
  \end{cases}
\end{equation}
where $s\!\cdot\!a$ means "the position after playing move $a$ in position $s$". This is
\textbf{minimax}: it assumes both sides play the best move available inside the part of the
tree we looked at.

\begin{worked}{minimax by hand}
White to move at the root, two plies deep, two moves each. Leaf evaluations (in pawns, from
White's point of view) are \emph{invented for the arithmetic}:

\begin{center}
\begin{tikzpicture}[level distance=13mm, sibling distance=22mm,
   every node/.style={font=\small},
   sq/.style={draw, rectangle, minimum width=8mm, minimum height=5mm, inner sep=1pt},
   leaf/.style={draw, rectangle, fill=black!4, minimum width=8mm, minimum height=5mm, inner sep=1pt}]
\node[sq, fill=BuildBlue!12] (root) {root, W}
  child {node[sq] (A) {$A$, B}
    child {node[leaf] {$+3$}}
    child {node[leaf] {$-1$}}}
  child {node[sq] (B) {$B$, B}
    child {node[leaf] {$+8$}}
    child {node[leaf] {$+2$}}};
\end{tikzpicture}
\end{center}

Work from the bottom up.
\begin{itemize}
\item Node $A$: \emph{Black} to move, so Black takes the minimum of $\{+3,-1\}$: $V(A) = -1$.
\item Node $B$: Black takes the minimum of $\{+8,+2\}$: $V(B) = +2$.
\item Root: \emph{White} to move, takes the maximum of $\{-1,+2\}$: $V(\text{root}) = +2$,
  achieved by playing $B$.
\end{itemize}

Notice what just happened at node $B$: the $+8$ leaf --- the most attractive number on the
board --- had \emph{no influence whatsoever}. Black simply would not go there. A very
common human error, and the one minimax is built to prevent, is letting a line's best-case
leaf colour your judgement of the line. The value of a move is what the opponent will
\emph{allow}, not what you hope for.
\end{worked}

\subsection{The negamax trick (you will need it in Chapter~\ref{ch:uct})}

Writing separate \texttt{max} and \texttt{min} cases is clumsy, and it becomes an actual bug
factory later when values flow up a Monte-Carlo tree. The standard fix is to score every
position \textbf{from the point of view of the side to move} instead of always from White's.
Then both players are maximisers, and the only thing that happens when you cross a ply
boundary is a change of sign:

\begin{equation}
  \label{eq:negamax}
  V(s) \;=\; \max_{a}\;\bigl[-\,V(s\!\cdot\!a)\bigr].
\end{equation}

Read it aloud: \emph{my value is the best I can do, where "good for me" is "bad for the
person who moves next".} Check it against the worked example: if $V(A) = -1$ from White's
view and it is Black to move at $A$, then from Black's own point of view $A$ is worth $+1$;
White's root value through $A$ is $-(+1) = -1$. Same answer, one rule.

\begin{pitfall}
This sign flip is responsible for a large fraction of all bugs in tree-search code, and for
a large fraction of all confusion when reading engine output. Leela's $Q$ values are
side-to-move-relative: a $Q$ of $+0.9$ in a node where Black is to move means \emph{Black}
is winning. Chapter~\ref{ch:byhand} makes you flip signs by hand thirty-odd times, on
purpose, until it stops being a thing you have to think about.
\end{pitfall}

\section{Alpha--beta: proving you do not need to look}

Minimax visits every leaf. Alpha--beta pruning gets the identical answer while skipping
large parts of the tree, using one observation: \emph{as soon as a line is proven bad enough
that the opponent will never allow it, stop analysing it}. You do this at the board already.

\begin{worked}{a cutoff, in chess words and in numbers}
White to move; two candidate moves, $A$ and $B$.
\begin{enumerate}
\item You analyse $A$ fully and conclude it gives $+2$. So you now have a guarantee: the
  root is worth \emph{at least} $+2$. Call that guarantee $\alpha = +2$.
\item You start on $B$ and look at Black's first reply $B_1$, which gives $+1$.
\item Stop. Black chooses at $B$, so $V(B) \le +1$. Black already has a way to hold $B$ to
  $+1$, which is worse for you than the $+2$ you have banked. Whatever Black's \emph{other}
  replies to $B$ do, they can only lower $V(B)$ further. So $B$ can never beat $A$, and the
  remaining replies to $B$ --- possibly millions of positions --- need never be examined.
\end{enumerate}
In chess: \emph{"I don't need to check the rest of this line; he has a good answer already."}
\end{worked}

\begin{established}
Alpha--beta returns exactly the same value as full minimax --- it prunes only branches that
provably cannot affect the root. With perfect move ordering it examines about $b^{d/2}$
leaves instead of $b^{d}$, i.e.\ the effective branching factor drops to $\sqrt{b}$
(Knuth \& Moore, 1975). With $b=31$, that is $b_{\text{eff}} \approx 5.6$: the difference
between depth 8 and depth 16 on the same hardware.
\end{established}

Two consequences worth carrying forward:

\begin{enumerate}
\item \textbf{Alpha--beta's power comes from move ordering.} The $\sqrt{b}$ result assumes
  you try the best move first. Classical engines therefore spend enormous effort on
  \emph{move-ordering heuristics} (try captures first, try the move that worked in a
  sibling node, etc.). Those heuristics are hand-written guesses about which move is good
  --- which is to say: \emph{classical engines already had a policy, they just wrote it by
  hand.} Leela's policy head is the learned version of exactly this. It is not a new kind of
  component; it is an old component that stopped being hand-written.
\item \textbf{Alpha--beta is all-or-nothing.} A branch is either searched or cut. There is
  no notion of "spend 3\% of my effort on this suspicious-looking move". Monte-Carlo search
  in Part~\textsc{ii} replaces that binary with a \emph{budget allocation}, and that
  difference is what makes Leela's search readable as a statement of interest --- the visit
  counts tell you how seriously each plan was taken.
\end{enumerate}

\section{The classical answer to Question 1: a hand-written formula}

Now Question~1. How did engines score a leaf before neural networks? By adding up features:

\begin{equation}
  \label{eq:handeval}
  f(s) \;=\; \sum_{i=1}^{k} w_i\,\phi_i(s),
\end{equation}
where each $\phi_i$ is a \textbf{feature} --- a countable property of the position --- and
$w_i$ is a hand-chosen \textbf{weight} saying how much that property is worth.

\begin{worked}{a complete (bad) evaluation function}
The simplest usable evaluation is material only, with the classical values.
\[
\begin{aligned}
\phi_1 &= (\text{White pawns}) - (\text{Black pawns}),   & w_1 &= 1.0\\
\phi_2 &= (\text{White knights}) - (\text{Black knights}),& w_2 &= 3.0\\
\phi_3 &= (\text{White bishops}) - (\text{Black bishops}),& w_3 &= 3.2\\
\phi_4 &= (\text{White rooks}) - (\text{Black rooks}),   & w_4 &= 5.0\\
\phi_5 &= (\text{White queens}) - (\text{Black queens}), & w_5 &= 9.0
\end{aligned}
\]
In a position where White is a knight up for two pawns:
$\phi_1 = -2$, $\phi_2 = +1$, rest zero, so
\[
f(s) = 1.0(-2) + 3.0(+1) = +1.0 .
\]
"White is a pawn better." A real classical evaluation adds dozens more terms --- passed
pawns, king safety, rook on an open file, bishop pair, mobility --- each with a weight
someone tuned. Stockfish's hand-written evaluation, before it went neural, had hundreds.
\end{worked}

Equation~\eqref{eq:handeval} is worth staring at, because \textbf{the neural network in
Part~\textsc{iv} is the same equation}. Really: the value head is $f(s) = \sum_i w_i
\phi_i(s)$ where the $\phi_i$ are computed by the layers below rather than written by a
person, and the $w_i$ are found by gradient descent rather than tuned by hand.
Chapter~\ref{ch:learned} makes that transition one step at a time. Nothing alien arrives.

\subsection{The two things a hand-written formula cannot do}

\begin{enumerate}
\item \textbf{It cannot represent what nobody thought to write down.} Every $\phi_i$ exists
  because a human noticed the pattern and coded a detector. "The dark squares around his
  king are permanently weak because the bishop that guarded them was traded and the pawn
  structure cannot repair it" is one feature to a strong player, and roughly no feature at
  all to a formula --- there is no counter for it.
\item \textbf{Its terms are added, so it cannot represent \emph{it depends}.} A passed pawn
  on d5 is worth a great deal in one structure and nothing in another; a knight on the rim
  is bad, except when it is heading for f5. Equation~\eqref{eq:handeval} scores each feature
  in isolation and sums. Interactions have to be hand-coded as yet more special-case terms,
  and the number of interactions explodes faster than anyone can write them.
\end{enumerate}

Compensation for sacrificed material is where both failures bite at once, which is exactly
your interest in this project: a sacrifice is precisely a position where the material term
is loud and wrong, and the compensating factors are diffuse, interacting, and unnamed.

\begin{keyidea}
Both of Leela's departures are attacks on these limits. The network learns the features
instead of being told them (fixing failure 1), and its layers multiply and gate information
rather than merely summing it (fixing failure 2). Chapter~\ref{ch:attention} is precisely
the story of how "it depends on what else is on the board" becomes an arithmetic operation.
\end{keyidea}

\section{What is actually in \texttt{engine/}}

To keep this concrete, here is the object this book is about, as it exists on your disk.

\begin{repobox}[title={in this repository}]
\begin{itemize}
\item \texttt{engine/lc0.exe} --- the search. A C++ program, version 0.32.1. It contains no
  chess knowledge beyond the rules: it builds and walks a tree
  (Chapters~\ref{ch:uct}--\ref{ch:engineering}).
\item \texttt{engine/BT3-768x15x24h-swa-2790000.pb.gz} --- a network. 15 transformer layers,
  768-dimensional square embeddings, 24 attention heads (Chapters~\ref{ch:attention},
  \ref{ch:heads}).
\item \texttt{engine/791556.pb.gz} --- a smaller, faster network. All the real numbers in
  this book come from this one, because it runs on this machine's CPU in seconds rather
  than minutes.
\item \texttt{engine/stockfish/} --- Stockfish 16.1, used in this book only as an
  independent referee for chess claims.
\end{itemize}
\end{repobox}

The split matters and is worth saying once plainly: \textbf{\texttt{lc0.exe} without a
network file cannot play chess at all}, and the network without \texttt{lc0.exe} plays about
as well as a strong club player glancing at a board for a tenth of a second. The engine you
know is the pair. Part~\textsc{ii} is the first half; Part~\textsc{iv} is the second;
Chapter~\ref{ch:byhand} is where you watch them talk to each other.

\begin{exercises}

\exhead[warm-up]{ch1:treesize} Using $b = 31.35$ (your corpus mean), how many positions are there at
depth 5? At depth 7? Roughly how many times bigger is depth 7 than depth 5, and why is that
ratio not a coincidence?

\exhead[the hardware fallacy]{ch1:hardware} A machine evaluates $2\times10^{7}$ positions per second and
searches to depth 9 in a certain time. You buy hardware $64\times$ faster. To what depth can
you now search in the same time? Show the calculation.

\exhead[branching factor in the endgame]{ch1:endgameb} Your corpus gives $b = 32.5$ in the middlegame and
$b = 15.8$ in the endgame. For a fixed budget of $10^{9}$ leaf evaluations, how deep can
exhaustive search go in each? Comment on why engine analysis "feels" so much more definite
in endgames.

\exhead[minimax]{ch1:minimax} White to move. Root has three children $A$, $B$, $C$ (Black to move at
each). Their children's leaf values, from White's point of view, are
$A: \{+1, +6, +2\}$, $B: \{+3, +3, +4\}$, $C: \{+9, -2, +7\}$.
Compute $V(A), V(B), V(C)$ and the root value, and give White's best move. Which leaf is the
largest number in the whole tree, and does it matter?

\exhead[negamax]{ch1:negamax} Redo the previous exercise using Equation~\eqref{eq:negamax}, i.e.\ with
every value expressed from the point of view of the player to move at that node. Write down
each node's value in that convention and confirm you get the same best move.

\exhead[alpha--beta by hand]{ch1:alphabeta} Take the tree from Exercise~\ref{ex:ch1:minimax}, searched
left to right ($A$ then $B$ then $C$; within each node, the leaves in the order listed).
After finishing $A$, what is $\alpha$? Walk through $B$ and $C$ and mark every leaf that
alpha--beta can skip. How many of the 9 leaves did you have to look at?

\exhead[ordering matters]{ch1:ordering} Same tree, but now search the children in the order $C$, $B$, $A$.
How many leaves are examined? What does the difference tell you about why classical engines
invest so heavily in move-ordering heuristics --- and what the learned equivalent of a
move-ordering heuristic is called in Leela?

\exhead[writing an evaluation]{ch1:evalfeature} Extend the material-only $f(s)$ of the worked example with
one positional feature of your own choosing (say, $\phi_6 =$ number of White rooks on open
files minus the same for Black). Pick a weight. Now describe, in one sentence each, two
different positions where your feature gives the wrong answer, and say which of the two
failure modes of §1.5.1 each one is.

\exhead[the compensation problem]{ch1:compensation} In one paragraph, explain why a piece sacrifice for a
long-term attack is the hardest case for Equation~\eqref{eq:handeval}, referring to both
failure modes. Keep the chess vague and the argument sharp --- this is a statement about the
\emph{form} of the equation, not about any particular sacrifice.

\exhead[looking ahead]{ch1:guess} Leela examines roughly 800--2{,}000 positions to choose a move in the
configuration this project uses, while a classical engine examines tens of millions. Before
reading on, write down your best guess at how that can possibly be competitive. Keep your
answer; you will check it against Chapter~\ref{ch:puct}.

\end{exercises}
